\section{Case II: The Instrument at the Wrong Layer}
\label{sec:wronglayer}

A two-block fast-weight model is trained on an associative binding
task ($K=32$ bindings per episode) against matched baselines, with two
registered readouts: a discrete task metric (episode-restricted recall
through the model's own continuation and LM head) and a continuous
probe metric $\rf$ (the fraction of targets a linear probe on a
tapped state recovers at cosine $\geq0.9$). In the third of three
consecutive rounds recorded as failures at $\rf=0.0$, the discrete
metric dissociated: recall accuracy $0.9990$ against chance $0.03125$
(baselines $0.0447$ and $0.0295$), the probe still exactly $0.0$ in
every arm.
% <!-- evidence: W1 -->
The dissociation, not the zero, is the finding (Rule 5).

\looseness=-1 An offline re-fit battery removed the probe-training explanations:
closed-form ridge regression, the global optimum for the linear probe
class, reproduced $\rf=0.0$ exactly, as did pinned SGD and MLP probes;
the trained probe's output was the analytically known
episode-membership direction (mean cosine $0.176$ against the
$1/\sqrt{32}\approx0.177$ ceiling; membership cosine $0.896$); the
ridge fit cleared its shuffled-tap control by at most $+0.059$
cosine.
% <!-- evidence: W2 -->
Meanwhile the full forward pass decodes the answer at $0.9957$
accuracy ($31.9\times$ chance) for the contender, both baselines at or
near chance: the information exists in the network and is not linearly
present at the tapped state.
% <!-- evidence: W3 -->

\looseness=-1 A state-zeroing experiment located the storage
(Figure~\ref{fig:taploc}): zeroing the
first block's state $S_0$ collapses recall to chance ($0.0286$);
zeroing $S_1$, the layer every deployed instrument tapped, leaves it
bit-unchanged at $0.9990$: the probed layer is causally inert for the
probed behavior. Placement alone is not the whole defect: a ridge
probe on the load-bearing $S_0$ \emph{also} reads $\rf=0.0$ (gaps over
shuffled controls $+0.002$ to $+0.063$ across all state-level taps),
and only the post-block-1, pre-LM-head hidden exposes the recall
linearly ($\rf=0.674$, gap $+0.800$), only in the arm that can perform
the task (the ablation's same tap reads $0.0$; Table~\ref{tab:taps}).
% <!-- evidence: W4 -->
The bindings are stored in a format no tested linear read on either
state exposes; the model's own downstream nonlinearity is the only
decoder found.

The same diagnosis exposed a second defect in passing: an auxiliary
classifier whose labels were, by construction, uniform given the
decoded quantity, so a perfect tap scores chance; its expected-null
had been read as confirmation for a round, hence Rule 3's
positive-control clause.
% <!-- evidence: W5 -->
